\documentclass[10pt,a4paper]{amsart}
\usepackage[margin=19mm]{geometry}
\usepackage{amsmath,amssymb,mathtools}
\usepackage[scr=rsfs,cal=euler]{mathalfa}
\usepackage{xfrac}
\usepackage[unicode,hidelinks,bookmarks=false]{hyperref}
\newcommand{\ie}{i.e.\ }
\newcommand{\cf}{cf.\ }
\newcommand{\resp}{respectively\ }
\newcommand{\Subsection}{\subsection}
\providecommand{\nobreakdash}{\nobreak\hbox{-}}
\setlength{\emergencystretch}{2em}
\multlinegap0pt
\usepackage[shortlabels]{enumitem}
\usepackage{amssymb}
\newtheorem{lemma}{Lemma}
\newcommand{\cS}{\mathcal{S}}
\title{On the Maximum Size of 2-Weakly Compatible Split Systems}
\author{Pei Wu, Stefan Gr\"unewald}
\date{Source: arXiv:2608.23275v1 (CC BY 4.0)}
\begin{document}
\maketitle

\noindent\emph{Source and licence.} On the Maximum Size of 2-Weakly Compatible Split Systems. Version arXiv:2608.23275v1 (CC BY 4.0), distributed by arXiv under the Creative Commons Attribution 4.0 International licence, http://creativecommons.org/licenses/by/4.0/, as recorded in the arXiv licence metadata for this exact version and shown on the abstract page https://arxiv.org/abs/2608.23275v1. Full licence notice: \texttt{licenses/arxiv-2608.23275-CC-BY-4.0.txt}.\par\smallskip

\noindent\textit{Editorial scope.} Counted unit: the article's lemma lem:nonnegative-charge (source0.tex lines 225--227). The article's complete original proof of it is delivered with it (source0.tex lines 229--303, one \texttt{proof} environment of 75 lines). No mathematics is written, repaired or re-derived: the statement and the proof are the article's own lines, in the article's own notation, and no retained line is substituted or re-flowed.  Retained uncounted as the source-local context the proof needs: lines 152--159.  Nothing was omitted from any retained span, so the package carries no omission record.  No retained line contains a citation, so the wrapper carries no bibliography and no bibliography entry.  No retained line contains a figure environment, an image-inclusion command or drawing code, so no image is delivered and the package declares no asset dependency.  Editorial formatting only, in addition: an AMS \texttt{amsart} preamble that restates the article's own declarations and macros used by the retained lines, namely the environments \texttt{lemma} on separate counters, together with the standard packages those lines need.  The retained environments print in this document as Lemma 1 in order of appearance, whereas the article prints the same items with the numbering of its own full document.  The complete native source of the article as distributed in the arXiv e-print for this exact version is bundled unchanged as \texttt{sources/208404/source0.tex}.
	\begin{lemma}\label{lem:hard-partners}
		For every graph $G$ representing a 2-weakly compatible split system, the following statements hold.
		\begin{enumerate}[label=(\roman*)]
			\item Every hard 6-vertex has at least one partner.
			\item If $A-a+b$ is a partner of a hard 6-vertex $A$, then $d^+(A-a)=d^-(A+b)=2$.
			\item No vertex is a partner of two distinct hard 6-vertices.
		\end{enumerate}
	\end{lemma}

	\begin{lemma}\label{lem:nonnegative-charge}
		Let $G=(V,E)$ be a digraph representing a 2-weakly compatible split system $\cS$ on $X$, and let $r\in X$ be the root of $G$. Then $C_2(A)\ge 0$ for every vertex $A$ with $5\le\#A\le n-5$.
	\end{lemma}

	\begin{proof}
		Suppose first that $C_1(A)\ge 0$. If $A$ sends no charge in the second discharging step, then $C_2(A)\ge C_1(A)\ge 0$. If $A$ sends charge, then $C_1(A)\ge 1$; moreover, Lemma~\ref{lem:hard-partners}(iii) implies that $A$ sends at most one unit. Thus $C_2(A)\ge 0$ in either case. It remains to show that every hard 6-vertex $A$ has a partner $A'$ with $C_1(A')\ge 1$.
		
		Let $A$ be a hard 6-vertex such that $A-a_i$ and $A+b_i$, for $i=1,2,3$, are its six neighbours and $A-a_i+b_i\in V$. Let $A-a_3+b_3$ be a partner of $A$. If $A$ has no other partner, then, without loss of generality, $A-a_1+b_2\in V$. Suppose that $d^+(A-a_3+b_3)=3$, and let
		\[
		A+b_3,\quad A-a_3+b_3+c_1,\quad A-a_3+b_3+c_2
		\]
		be the out-neighbours of $A-a_3+b_3$. If $b_2\notin\{c_1,c_2\}$, then the vertices
		\[
		A-a_3+b_3+c_1,\quad A-a_3+b_3+c_2,\quad A-a_2+b_2,\quad A-a_1+b_2
		\]
		contradict 2-weak compatibility. If $b_2\in\{c_1,c_2\}$, say $b_2=c_2$, then the vertices
		\[
		A-a_2+b_2,\quad A-a_1+b_2,\quad A,\quad A-a_3+b_3+b_2
		\]
		contradict 2-weak compatibility, because the taxa $a_1,a_2,a_3,b_2,b_1,r$ induce the forbidden configuration. Hence $d^+(A-a_3+b_3)\le 2$, and, by symmetry, $d^-(A-a_3+b_3)\le 2$. Thus $C_1(A-a_3+b_3)\ge 1$, and the second discharging rule gives $C_2(A)=0$.
		
		It remains to consider the case in which no vertex $A-a_i+b_j$ with $i\ne j$ and $\{i,j\}\subseteq\{1,2,3\}$ belongs to $V$. Then all three vertices $A-a_i+b_i$, for $i=1,2,3$, are partners of $A$. The assertion holds if one of these vertices has degree at most 4. We may therefore assume that every partner of $A$ has degree at least 5. It follows that two partners of $A$ have out-degree 3 or two have in-degree 3. Without loss of generality, assume that
		\[
		d^+(A-a_1+b_1)=d^+(A-a_2+b_2)=3.
		\]
		For $i=1,2$, let
		\[
		A-a_i+b_i+c_i^1,\quad A-a_i+b_i+c_i^2
		\]
		be the two out-neighbours of $A-a_i+b_i$ other than $A+b_i$. We have
		\[
		b_3\notin \{c_1^1,c_1^2\}\cap\{c_2^1,c_2^2\},
		\]
		because otherwise the vertices
		\[
		A,\quad A+b_3-a_3,\quad A+b_3+b_2-a_2,\quad A+b_3+b_1-a_1
		\]
		would not be 2-weakly compatible: the taxa $a_1,a_2,a_3,b_3$, together with two elements of $X-A-b_1-b_2-b_3$, would induce the forbidden configuration. Hence, without loss of generality, $b_3\notin\{c_1^1,c_1^2\}$.
		
		We have $A+b_3-a_3-a_1\notin V$, because this vertex together with
		\[
		A,\quad A+b_1-a_1+c_1^1,\quad A+b_1-a_1+c_1^2
		\]
		would violate 2-weak compatibility. Moreover, $d^-(A-a_3+b_3)=3$ is impossible; otherwise, the two in-neighbours of $A-a_3+b_3$ other than $A-a_3$, together with $A+b_1-a_1+c_1^1$ and $A+b_1-a_1+c_1^2$, would not be 2-weakly compatible.
		
		We have shown that
		\[
		d^+(A-a_1+b_1)=d^+(A-a_2+b_2)=3
		\]
		implies $d^-(A-a_3+b_3)=2$. By symmetry, and because $d(A-a_i+b_i)\ge 5$ for $i=1,2,3$, we may assume that
		\[
		d^+(A-a_i+b_i)=3\quad\text{and}\quad d^-(A-a_i+b_i)=2
		\]
		for $i=1,2,3$. Using the same argument that allowed us to assume $b_3\notin\{c_1^1,c_1^2\}$, we may also assume that $b_1\notin\{c_2^1,c_2^2\}$. Let $A-a_1+b_1-d_1$ and $A-a_1$ be the in-neighbours of $A-a_1+b_1$. We cannot have
		\[
		\{A-a_1+b_1-d_1+c_1^1,\ A-a_1+b_1-d_1+c_1^2\}\subseteq V,
		\]
		because these two vertices, together with either $A-a_2,A-a_3$ if $d_1\notin\{a_2,a_3\}$, or with $A-a_1,A-d_1$ otherwise, would not be 2-weakly compatible. Assume, without loss of generality, that
		\[
		A-a_1+b_1-d_1+c_1^2\notin V.
		\]
		If $d^-(A-a_1+b_1+c_1^2)>1$, choose $f\in X\setminus\{c_1^2,d_1\}$ such that
		\[
		A-a_1+b_1+c_1^2-f\in V.
		\]
		Then $f\ne b_1$, because otherwise the vertices $A-a_1+c_1^2,A+b_1,A+b_2,A+b_3$ are not 2-weakly compatible if $c_1^2\notin\{b_2,b_3\}$, while $A-a_1+b_1$ is not a partner of $A$ otherwise. Furthermore, $a_2\notin\{d_1,f\}$, because otherwise the vertices
		\[
		A-a_2+b_2+c_2^1,\quad A-a_2+b_2+c_2^2,\quad A+b_2,
		\]
		together with $A-a_1-a_2+b_1$ if $d_1=a_2$, or with $A-a_1-a_2+c_1^2$ if $f=a_2$, would not be 2-weakly compatible. However, $f\ne b_1$, $a_2\notin\{d_1,f\}$, and $b_1\notin\{c_2^1,c_2^2\}$ imply that the vertices
		\[
		A-a_1+b_1+c_1^2-f,\quad A-a_1+b_1-d_1,\quad A-a_2+b_2+c_2^1,\quad A-a_2+b_2+c_2^2
		\]
		are not 2-weakly compatible, because the taxa $a_2,b_1,c_2^1,c_2^2,d_1,f$ induce the forbidden configuration. Therefore,
		\[
		d^-(A-a_1+b_1+c_1^2)=1.
		\]
		Thus $A-a_1+b_1+c_1^2$ sends one unit of charge to $A-a_1+b_1$ in the first discharging step, and $A-a_1+b_1$ sends one unit of charge to $A$ in the second step. Hence $C_2(A)\ge 0$.
	\end{proof}

\end{document}
